\documentclass[11pt]{scrartcl}
\usepackage{ebb}



\title {Trigonometria}
\subtitle {Entrenamientos Nacionales Centros Pagmos}

\author{Emmanuel Buenrostro}

\begin{document}
\maketitle

\section{Problemas}
\epigraph{The cold water does not get warmer if you jump late}{ }

\textit{Nota:} Agradecimientos a Mateo por dejarme copiar problemas de su lista.
\subsection{Nivel I}
\problemdiff{0}

\begin{problem} %[AHSME 1963/34]
    En $\triangle ABC$, el lado $a=\sqrt{3}$, el lado $b=\sqrt{3}$, y el lado $c>3$. Sea $x$ el mayor numero tal que la magnitud, en grados, del angulo opuesto al lado $c$ excede $x$. Entonces $x$ es igual a:
\end{problem}



\problemdiff{3}

\begin{problem} %[ASHME 1991/29]
    El triangulo equilatero $ABC$ tiene $P$ sobre $AB$ y $Q$ sobre $AC$. El triangulo se dobla a lo largo de $PQ$ de modo que el vertice $A$ ahora queda en $A'$ sobre el lado $BC$. Si $BA'=1$ y $A'C=2$, calcula la longitud de  $PQ$
\end{problem}

\begin{problem} %[AIME I 2012/12]
Sea $\triangle ABC$ un triangulo rectangulo con angulo recto en $C$. Sean $D$ y $E$ puntos sobre $\overline{AB}$ con $D$ entre $A$ y $E$ tales que $\overline{CD}$ y $\overline{CE}$ trisecan $\angle C$. Si $\frac{DE}{BE}=\frac{8}{15}$, Encuentra $\tan B$.
\end{problem}

\begin{problem} %[AIME II 2016/10]
    El triangulo $ABC$ esta inscrito en el circulo $\omega$. Los puntos $P$ y $Q$ estan sobre el lado $\overline{AB}$ con $AP<AQ$. Los rayos $CP$ y $CQ$ intersectan nuevamente a $\omega$ en $S$ y $T$ (distintos de $C$), respectivamente. Si $AP=4$, $PQ=3$, $QB=6$, $BT=5$, y $AS=7$, calcula $ST$.
\end{problem}
\subsection{Nivel II}
\begin{problem}%[``Lema de Irán'']

	El incírculo de $\triangle ABC$ con centro en $I$ es tangente a $\ol{BC}$, $\ol{CA}$, $\ol{AB}$ en $D$, $E$, $F$, respectivamente. Sean $M$
	y $N$ los puntos medios de $\ol{BC}$ y $\ol{AC}$, respectivamente. Demuestra que las líneas $MN$, $EF$ y $BI$ concurren.

\end{problem}

\begin{problem}
    

	En la configuración del problema anterior, demuestra que $DI$ y $EF$ concurren con la mediana de $\triangle ABC$ desde $A$. 


\end{problem}

\begin{problem}
    

	Sea $ABC$ un triángulo con circuncírculo $\omega$ y sea $T$ la intersección de las tangentes a $\omega$ por $B$ y $C$. Sea $M$ el punto medio
	de $\ol{BC}$. Demuestra que $\angle BAT = \angle MAC$.

\end{problem}



\problemdiff{6}
\begin{problem}%[Japón 2003 P1]

	El punto $P$ está dentro del triángulo acutángulo $ABC$. La línea $BP$ intersecta a $AC$ en $E$ y la línea $CP$ intersecta a $AB$ en $F$.
	Dado que $PE = CE$ y que $AF = FB = CP$, encuentra la medida de $\angle BFC$.

\end{problem}

\begin{problem} %[OMM 2023 P3]

	Sea $ABCD$ un cuadrilátero convexo. $M$, $N$, $K$ son los puntos medios de los segmentos $AB$, $BC$ y $CD$, respectivamente. Existe un 
	punto $P$ dentro del cuadrilátero tal que $\angle BPN = \angle PAD$ y $\angle CPN = \angle PDA$. Demuestra que $AB \cdot CD = 4PM \cdot PK$.

\end{problem}
\begin{problem}%[IMO 2009 P4]

    Sea $ABC$ un triángulo con $AB=AC$. Las bisectrices de los ángulos $\angle CAB$ y $\angle ABC$ intersecan los lados $BC$ y $CA$ en $D$ y $E$, respectivamente. 
	Sea $K$ el incentro del triángulo $ADC$. Si $\angle BEK=45^\circ$, encuentra todos los valores posibles del ángulo $\angle CAB$.
    
\end{problem}

\subsection{Nivel III}
\problemdiff{8}
\begin{problem} %[IMO 2022/4]
Sea $ABCDE$ un pentagono convexo tal que $BC=DE$. Supongamos que existe un punto $T$ dentro de $ABCDE$ con $TB=TD$, $TC=TE$ y $\angle ABT=\angle TEA$. Sea la recta $AB$ que intersecta a las rectas $CD$ y $CT$ en los puntos $P$ y $Q$, respectivamente. Supongamos que los puntos $P,B,A,Q$ aparecen sobre su recta en ese orden. Sea la recta $AE$ que intersecta a $CD$ y $DT$ en los puntos $R$ y $S$, respectivamente. Supongamos que los puntos $R,E,A,S$ aparecen sobre su recta en ese orden. Demuestra que los puntos $P,S,Q,R$ yacen sobre una misma circunferencia.
\end{problem}
\begin{problem} %[IMOSL 2017/G1]
    Sea $ABCDE$ un pentagono convexo tal que $AB=BC=CD$, $\angle EAB=\angle BCD$, y $\angle EDC=\angle CBA$. Demuestra que la recta perpendicular desde $E$ a $BC$ y los segmentos $AC$ y $BD$ son concurrentes.
\end{problem}

\begin{problem} %[IMOSL 2012/G2]
    Sea $ABCD$ un cuadrilatero ciclico cuyas diagonales $AC$ y $BD$ se intersectan en $E$. Las extensiones de los lados $AD$ y $BC$ mas alla de $A$ y $B$ se intersectan en $F$. Sea $G$ el punto tal que $ECGD$ es un paralelogramo, y sea $H$ la imagen de $E$ bajo la reflexion respecto a $AD$. Demuestra que $D,H,F,G$ son conciclicos.
\end{problem}

\begin{problem} %[IMO 2025/2]
    Sean $\Omega$ y $\Gamma$ circunferencias de centros $M$ y $N$, respectivamente, tales que el radio
	de $\Omega$ es menor que el radio de $\Gamma$. Supongamos que las circunferencias $\Omega$ y $\Gamma$
	se cortan en dos puntos distintos $A$ y $B$. La recta $MN$ corta a $\Omega$ en $C$ y a $\Gamma$ en $D$, de forma
	que los puntos $C$, $M$, $N$ y D están sobre esa recta en ese orden. Sea $P$ el circuncentro del triángulo $ACD$.
	La recta $AP$ corta de nuevo a $\Omega$ en $E \neq A$. La recta $AP$ corta de nuevo a $\Gamma$ en $F \neq A$. Sea $H$ el
	ortocentro del triángulo $PMN$.
	
	Demuestre que la recta paralela a $AP$ que pasa por $H$ es tangente al circuncículo del triángulo $BEF$.
	
\end{problem}

\end{document}